The modelling of physical systems has traditionally relied on mathematical formulations derived from fundamental laws, followed by numerical methods to solve the resulting differential equations \cite{strogatz2015nonlinear}. These approaches have been highly successful for a wide range of systems; however, obtaining accurate solutions can become increasingly challenging when the governing dynamics are nonlinear, high-dimensional, or sensitive to initial conditions \cite{strogatz2015nonlinear}. In recent years, machine learning techniques have emerged as an alternative approach, allowing complex relationships between variables to be approximated directly from data \cite{brunton2019data}.

Artificial neural networks have demonstrated significant potential in scientific computing due to their ability to approximate nonlinear functions \cite{hornik1989multilayer}. The universal approximation theorem states that feedforward neural networks can represent continuous functions with arbitrary accuracy given sufficient model capacity \cite{hornik1989multilayer}. However, this theoretical property does not guarantee efficient training or accurate generalization, particularly when the underlying dynamics exhibit complex temporal behaviour \cite{goodfellow2016deep}. Nevertheless, neural networks provide a promising framework for learning the evolution of dynamical systems directly from generated data.

The performance of neural networks depends strongly on the characteristics of the dynamics being approximated. Linear systems generally produce smooth and predictable trajectories, whereas nonlinear systems can introduce phenomena such as amplitude-dependent frequencies, complex oscillations, and chaotic behaviour \cite{strogatz2015nonlinear}. Systems exhibiting sensitivity to initial conditions present an additional challenge, since small variations in the initial state can lead to significantly different trajectories over time \cite{lorenz1963deterministic}. Understanding how these properties influence prediction accuracy is essential for evaluating the limitations of data-driven approaches in physical modelling.

In this work, the capability of a multilayer perceptron (MLP) to approximate several deterministic dynamical systems with increasing levels of complexity is investigated. Six systems are analysed: the harmonic oscillator, damped harmonic oscillator, forced harmonic oscillator, Duffing oscillator, simple pendulum, and double pendulum. A common neural network architecture is used throughout the experiments, and systematic parameter sweeps are performed to study how variations in physical parameters influence prediction performance.

The main objective of this work is to determine whether prediction difficulty is primarily related to the mathematical nature of the governing equations or to the complexity of the resulting trajectories. By comparing linear, nonlinear, and chaotic systems under a consistent framework, this study provides insight into the capabilities and limitations of simple feedforward neural networks for approximating deterministic physical dynamics.